\emph{suppose that for every point @@M0@@ and every irreducible component @@M1@@ of @@M2@@, the irreducible components of @@M3@@ have dimension @@M4@@, (which is the case if the irreducible components of the fibers of @@M5@@ have dimension @@M6@@). Then the induced morphism between the Picard functors}
@@M7@@
\emph{is of finite type.}

\emph{Moreover, if @@M8@@ is integral, there exists a nonempty open subset @@M9@@ of @@M10@@ such that the Picard functors restricted to @@M11@@ are representable and the induced morphism}
@@M12@@
\emph{is quasi-affine (of finite type).}

Let us first prove the first assertion. As in the proof of 3.6, one sees by noetherian induction that one may replace @@M13@@ by any nonempty open subset, and one easily reduces to the case where @@M14@@ is integral and @@M15@@ is integral over @@M16@@, always replacing @@M17@@ by @@M18@@ when @@M19@@ is replaced by @@M20@@\,; this done, one reduces in the same way to the case where @@M21@@ is normal over @@M22@@, using the following lemma in place of XII 1.3.

\sgaStableTarget{sga6.e13.lem.3-9}\textbf{Lemma 3.9.} \emph{Let @@M23@@ be an integral noetherian scheme and @@M24@@ an @@M25@@-scheme, integral over @@M26@@. Then there exist a nonempty open subset @@M27@@ of @@M28@@, an integral scheme @@M29@@ finite and faithfully flat over @@M30@@, an @@M31@@-scheme @@M32@@ normal and integral over @@M33@@, and a finite surjective @@M34@@-morphism @@M35@@.}

